
\documentclass{article}
\usepackage{colortbl}
\usepackage{makecell}
\usepackage{multirow}
\usepackage{supertabular}

\begin{document}

\newcounter{utterance}

\centering \large Interaction Transcript for game `ifeval', experiment `ifeval', episode 59 with reflexion{-}Qwen3.5{-}9B.
\vspace{24pt}

{ \footnotesize  \setcounter{utterance}{1}
\setlength{\tabcolsep}{0pt}
\begin{supertabular}{c@{$\;$}|p{.15\linewidth}@{}p{.15\linewidth}p{.15\linewidth}p{.15\linewidth}p{.15\linewidth}p{.15\linewidth}}
    \# & \multicolumn{2}{c}{Player} && \multicolumn{2}{c}{Game Master} \\
    \hline

    \theutterance \stepcounter{utterance}  
    & & & \multicolumn{4}{p{0.6\linewidth}}{
        \cellcolor[rgb]{0.9,0.9,0.9}{
            \makecell[{{p{\linewidth}}}]{
                \texttt{\tiny{[P1$\langle$GM]}}
                \texttt{You have attempted this same task before and did not succeed. These are the plans you wrote after those attempts.} \\
\texttt{Trial \#1: New plan: I failed because my response contained multiple commas and the letter "c" in words like "aims," "study," "patterns," and "decades." Additionally, my word count was likely insufficient. Next time, I will write the entire proposal as a single long sentence or a series of sentences strictly avoiding commas. I will carefully scan every word to ensure no "c" appears, replacing words like "research" with "study," "complex" with "hard," and "access" with "entry." I will also expand on the methodology and expected outcomes using simple vocabulary to guarantee the text reaches the 250{-}word minimum while adhering to the negative constraints.} \\
\\ 
\\ 
\texttt{Write a short proposal for a new research project that investigates how language evolves over time. I want to make it challenging, so:} \\
\texttt{1. Do not include any commas in your response.} \\
\texttt{2. Do not include the letter "c" anywhere in your response.} \\
\texttt{3. Your response should contain at least 250 words.} \\
            }
        }
    }
     \\ \\

    \theutterance \stepcounter{utterance}  
    & \multicolumn{4}{p{0.6\linewidth}}{
        \cellcolor[rgb]{0.561,0.737,0.561}{
            \makecell[{{p{\linewidth}}}]{
                \texttt{\tiny{[P1$\rangle$GM]}}
                \texttt{I propose a new study to track how speech changes across many years. This project will look at old books and new posts to see how words shift in meaning. We will pick common terms and check their use in different places. The team will read thousands of pages from the past one hundred years. They will also check modern social media feeds for the same terms. This method helps us see fast changes that happen today. We will count how often each word appears in each time period. The data will show clear patterns of growth or decline. Some words might gain new meanings while others lose old ones. We will map these shifts on a big chart for everyone to see. The results will help teachers understand how language grows naturally. Students can learn that words are not fixed but always moving. This work avoids hard theories and uses simple counting instead. We will share our findings in a clear report for all readers. The goal is to show how people create new ways to speak. Every day we find new phrases that did not exist before. Our team will watch these new phrases become common over time. We will also look at why some changes stick while others fade away. This study will run for three full years to get enough data. We will publish the final results in a journal for scholars to read. Many people will benefit from seeing how language lives and breathes. This project proves that studying words is fun and very useful. We will start this work soon to catch the latest changes. The future of speech depends on what we do now today. Our team is ready to begin this important and exciting journey. We will use only free tools to keep costs very low. Everyone can join this effort to learn about our shared history. The world changes fast and so does the way we talk. This study captures that speed in a simple and clear way. We hope to inspire others to join this great adventure soon.} \\
            }
        }
    }
    & & \\ \\

    \theutterance \stepcounter{utterance}  
    & & & \multicolumn{2}{p{0.3\linewidth}}{
        \cellcolor[rgb]{0.9,0.9,0.9}{
            \makecell[{{p{\linewidth}}}]{
                \texttt{\tiny{[GM$|$GM]}}
                {'punctuation:no_comma': {}, 'keywords:letter_frequency': {'let_relation': 'less than', 'letter': 'c', 'let_frequency': 1}, 'length_constraints:number_words': {'relation': 'at least', 'num_words': 250}}
            }
        }
    }
    & & \\ \\

    \theutterance \stepcounter{utterance}  
    & & & \multicolumn{2}{p{0.3\linewidth}}{
        \cellcolor[rgb]{0.9,0.9,0.9}{
            \makecell[{{p{\linewidth}}}]{
                \texttt{\tiny{[GM$|$GM]}}
                \texttt{game\_result = LOSE} \\
            }
        }
    }
    & & \\ \\

    \theutterance \stepcounter{utterance}  
    & & \multicolumn{4}{p{0.6\linewidth}}{
        \cellcolor[rgb]{0.95,0.95,0.95}{
            \makecell[{{p{\linewidth}}}]{
                \texttt{\tiny{[P1$\leftrightarrow$P]}}
                \texttt{Post{-}episode Reflexion (not sent to GM):} \\
\texttt{New plan: I failed because my response contained multiple commas and the letter "c" in words like "aims," "study," "patterns," and "decades." Additionally, my word count was likely insufficient. Next time, I will write the entire proposal as a single long sentence or a series of sentences strictly avoiding commas. I will carefully scan every word to ensure no "c" appears, replacing words like "research" with "study," "complex" with "hard," and "access" with "entry." I will also expand on the methodology and expected outcomes using simple vocabulary to guarantee the text reaches the 250{-}word minimum while adhering to the negative constraints.} \\
            }
        }
    }
    & \\ \\

\end{supertabular}
}

\end{document}
